\section{Stein's square function}\label{sec:stein-square}

Stein introduced a square function to study the convergence of
Bochner--Riesz means \cite[\S18, (18.1)]{Stein1958Square}.
For the Euclidean means \(B_r^\alpha\), with multiplier
\((1-|\xi|^2/r^2)_+^\alpha\), it measures the change with radius:
\[
 \mathcal G_\alpha f(x)=
 \left(\int_0^\infty |r\partial_r B_r^\alpha f(x)|^2
                         \frac{\dd r}{r}\right)^{1/2}.
\]
Carbery proved the sharp bounds in the plane
\cite[pp.~409--416]{Carbery1983Square}. In dimensions \(n\ge3\),
Lee, Rogers and Seeger proved them for \(p\ge 2+4/n\)
\cite[Theorem~1.1]{LRS2012Square}. Gan, Jing and Wu reached
\(p\ge13/4\) in dimension three \cite[Theorem~1.1]{GJW2024Square},
and Gan, Oh and Wu further enlarged the range in dimensions \(n\ge4\)
\cite[Theorem~1.1]{GOW2025Square}. Here the wave estimate
\eqref{eq:cs-unit-wave} gives the full range \(2\le p<\infty\).

\begin{theorem}\label{thm:stein-square}
Let \(n\ge2\) and \(2\le p<\infty\). If
\[
 \alpha>\max\left\{\tfrac12,\,n\left(\tfrac12-\tfrac1p\right)\right\},
\]
then \(\|\mathcal G_\alpha f\|_p\le C_{n,p,\alpha}\|f\|_p\) for
\(f\in\mathcal S(\RR^n)\), and \(\mathcal G_\alpha\) extends to \(L^p(\RR^n)\).
\end{theorem}

The threshold can be read off from one shell. At distance \(\delta\)
from \(|\xi|=s\), differentiating the mean gives a factor
\(\delta^{\alpha-1}\). The square norm in \(s\) gives back
\(\delta^{1/2}\), while the wave estimate at times of order
\(\delta^{-1}\) costs \(\delta^{-s_+(p)-\varepsilon}\). The resulting
powers sum if \(\alpha>1/2+s_+(p)\). We first make this calculation
on \(1\le s\le2\). We then combine the dyadic intervals of radii inside
the square norm, before taking the spatial \(L^p\) norm.

\begin{lemma}[One interval of radii]\label{lem:square-shell}
Let \(2<p<\infty\), \(\varepsilon>0\), \(\sigma=s_+(p)+\varepsilon\),
and \(H=L^2([1,2],\dd s)\). Fix \(A\in C_c^\infty(\RR\times(0,\infty))\)
and \(\varphi\in C_c^\infty(\RR)\). For \(0<\delta\le1\), let
\(T_\delta f(x,s)\) be the spatial Fourier multiplier
\[
 A(s,|\xi|)\,\varphi\bigl((s-|\xi|)/\delta\bigr)
\]
applied to \(f\). Then
\[
 \|T_\delta f\|_{L^p_x(H)}\lesssim \delta^{1/2-\sigma}\|f\|_p,
 \qquad
 \|T_\delta f\|_{L^2_x(H)}\lesssim \delta^{1/2}\|f\|_2,
\]
with constants independent of \(\delta\).
\end{lemma}

\begin{proof}
Temporarily omit \(A\), allow \(s\in\RR\), and denote the resulting
operator by \(V_\delta\). For \(g\) supported in a fixed compact
annulus in frequency, Plancherel in \(s\) gives, at each \(x\),
\[
 \|V_\delta g(x,\cdot)\|_2
   =(2\pi)^{-1/2}\delta
          \|\widehat\varphi(\delta t) U(t)g(x)\|_{L^2_t}.
\]
On \(R\le |t|\le 2R\), \(R\ge1\), H\"older in \(t\) and the rescaled
wave estimate \eqref{eq:cs-unit-wave} give
\[
 \|U(t)g\|_{L^p_xL^2_t}
 \lesssim R^{1/2-1/p}R^{1/p+\sigma}\|g\|_p.
\]
On \(|t|\le1\) the annular wave kernels have uniformly bounded
\(L^1\) norms. For \(b>1/2+\sigma\), summing over \(R=1,2,4,\ldots\)
now gives
\[
 \|V_\delta g\|_{L^p_x L^2_s(\RR)}
 \lesssim \delta\left(1+\sum_R R^{1/2+\sigma}(1+\delta R)^{-b}\right)
             \|g\|_p
 \lesssim \delta^{1/2-\sigma}\|g\|_p .
\]
To restore \(A\), Fourier invert it in \(s\). The multipliers
\(\widehat A_1(\eta,|D|)\), where the hat denotes transformation in
the first argument, are supported in a fixed annulus and have \(L^p\)
operator norms \(O_b((1+|\eta|)^{-b})\) for every \(b>0\), by smoothness
and compact support. Integrating
\(e^{\mathrm is\eta}V_\delta(\widehat A_1(\eta,|D|)f)\) in \(\eta\)
with measure \(\dd\eta/(2\pi)\) proves the first bound. For the second,
spatial Plancherel suffices: at each frequency the multiplier is bounded
and is supported on an \(s\)-interval of length \(O(\delta)\).
\end{proof}

To combine the intervals of radii, we use Seeger's method for passing
from bounds on individual frequency annuli to a multiplier estimate
\cite[Theorem~1]{Seeger1988Dilates}. Jones, Seeger and Wright stated
a square function version for dilates of annular multipliers
\cite[p.~6737]{JSW2008Dilates}; Guo, Roos and Yung proved it
\cite[Proposition~4.2 and \S7]{GRY2020Dilates}. Here we allow a different
input at each scale and assume decay of the Hilbert-valued kernel and
its first derivative. On each test cube only \(O(1+\log(1/\delta))\)
scales need the \(L^p\) bound. H\"older over those scales gives the
logarithmic power \(1/2-1/p\), which the strict inequality in
Theorem~\ref{thm:stein-square} absorbs.

\begin{lemma}[Combining dilates]\label{lem:square-dilates}
Let \(H\) be a Hilbert space, \(2<p<\infty\), \(0<\delta\le1\), and
\(N=n+3\). Suppose \(K:\RR^n\to H\) is continuously differentiable and
\[
 \|K(x)\|_H+\|\nabla K(x)\|_{H^n}
             \le C_*\delta^{-N}(1+|x|)^{-N}.
\]
Write \(S^{(k)}f=2^{nk}K(2^k\cdot)*f\), and let \(B_q\) bound the norm
of \(K*\) from \(L^q\) to \(L^q(H)\), \(q=2,p\). For every finite
sequence of Schwartz functions \(f_k\),
\[
 \left\|\left(\sum_k\|S^{(k)}f_k\|_H^2\right)^{1/2}\right\|_p
 \lesssim_{n,p}
 \left(B_2+C_*\delta+[1+\log(1/\delta)]^{1/2-1/p}B_p\right)
 \left\|\left(\sum_k|f_k|^2\right)^{1/2}\right\|_p .
\]
\end{lemma}

\begin{proof}
Put \(b_k=\|S^{(k)}f_k\|_H\),
\(h=(\sum_k b_k^2)^{1/2}\), and \(W=(\sum_k|f_k|^2)^{1/2}\).
Let \(M\) denote the Hardy--Littlewood maximal operator over cubes and
\(L=\lceil(N+2)\log_2(2/\delta)\rceil\).
On a dyadic cube \(Q\) of side \(2^{-j}\) and center \(y\), the
terms with \(k\le j-L\) are nearly constant. We compare \(h\)
with \(c_Q=(\sum_{k\le j-L}b_k(y)^2)^{1/2}\).
For \(x,z\in Q\) and \(k\le j-L\), the derivative bound gives
\[
 \|S^{(k)}f_k(x)-S^{(k)}f_k(y)\|_H
       \lesssim C_*\delta^{-N} 2^{k-j}MW(z).
\]
Indeed \(2^k|x-z|\lesssim1\), and convolution with
\(2^{nk}(1+2^k|\cdot|)^{-N}\) is bounded by a constant times \(M\).
Summing the squares and taking their square root costs at most
\(C C_*\delta MW(z)\).

For \(k\ge j+L\), split \(f_k\) on \(5Q\) and its complement. The
local terms satisfy
\[
 \frac1{|Q|}\int_Q\sum_{k\ge j+L}
       \|S^{(k)}(f_k\mathbf1_{5Q})(x)\|_H^2\dd x
       \lesssim B_2^2 M(W^2)(z),\qquad z\in Q,
\]
because dilation preserves the \(L^2\) operator norm. For the terms
outside \(5Q\), summing the kernel bound on successive doubled cubes
gives, for \(x,z\in Q\),
\[
 \|S^{(k)}(f_k\mathbf1_{(5Q)^c})(x)\|_H
       \lesssim C_*\delta^{-N}2^{-(N-n)(k-j)}MW(z).
\]
Their square sum again has square root at most \(C C_*\delta MW(z)\).

There remain fewer than \(2L\) indices. H\"older in this finite set
bounds \(\bigl(|Q|^{-1}\int_Q\sum_{|k-j|<L} b_k^2\bigr)^{1/2}\) by
\((2L)^{1/2-1/p}F(z)\), where
\(F=(\sum_k [M(b_k^2)]^{p/2})^{1/p}\). Since \(p/2>1\), the scalar
maximal theorem and the scale invariance of \(B_p\) give
\[
 \|F\|_p^p\lesssim\sum_k\|b_k\|_p^p
       \le B_p^p\sum_k\|f_k\|_p^p\le B_p^p\|W\|_p^p.
\]
By the triangle inequality for the square norm and Cauchy--Schwarz on
\(Q\), the mean of \(|h-c_Q|\) is thus at most a constant times
\(
 C_*\delta MW(z)+B_2 M(W^2)(z)^{1/2}+L^{1/2-1/p}F(z).
\)
Replacing \(c_Q\) by the mean \(h_Q\) costs at most a factor of two.
The dyadic sharp function inequality
\cite[Chapter~IV, \S3.6, Corollary~1]{Stein1993}
and the maximal theorem on \(L^p\) and \(L^{p/2}\) now prove the
assertion. The hypothesis of that inequality at a smaller exponent holds:
$\|h\|_2^2=\sum_k\|S^{(k)}f_k\|_{L^2(H)}^2
\le B_2^2\sum_k\|f_k\|_2^2<\infty$, and $2<p$.
\end{proof}

\begin{proof}[Proof of Theorem~\ref{thm:stein-square}]
For \(p=2\), spatial Plancherel and \(u=|\xi|^2/r^2\) give the identity
\[
 \|\mathcal G_\alpha f\|_2^2
 =2\alpha^2\int_0^1 u(1-u)^{2\alpha-2}\dd u\,\|f\|_2^2
 =\frac{\alpha}{2\alpha-1}\|f\|_2^2.
\]
Suppose \(p>2\). Choose \(\varepsilon>0\) so that
\(\sigma=s_+(p)+\varepsilon<\alpha-1/2\), possible since
\(1/2+s_+(p)=\max\{1/2,n(1/2-1/p)\}\).
Write \(r=2^ks\), \(1\le s\le2\), and \(\rho=2^{-k}|\xi|\).
On this interval \(\dd r/r=\dd s/s\) is comparable to \(\dd s\).
The differentiated multiplier equals
\[
 m(s,\rho)=2\alpha\rho^2s^{-2\alpha}(s+\rho)^{\alpha-1}
             (s-\rho)^{\alpha-1}\mathbf1_{\{\rho<s\}} ,
\]
where the last power is taken only for \(\rho<s\).
Near the boundary \(\rho=s\), split \(s-\rho>0\) into dyadic intervals
of size \(\delta=2^{-j}\), \(j\ge0\). Each piece is
\(\delta^{\alpha-1}A(s,\rho)\varphi((s-\rho)/\delta)\):
the factors preceding \((s-\rho)^{\alpha-1}\) times a smooth boundary
cutoff give a fixed amplitude \(A\) supported away from \(\rho=0\);
the power on a unit shell is absorbed in
\(\varphi\in C_c^\infty((1/2,2))\).

For fixed \(\delta\), Lemma~\ref{lem:square-shell} bounds the one-interval
norms of \(T_\delta\). Its kernel and first spatial derivatives satisfy
the hypothesis of Lemma~\ref{lem:square-dilates} with \(C_*\) independent
of \(\delta\): \(N\) integrations by parts in frequency cost at most
\(\delta^{-N}\) on a fixed compact annulus, uniformly in \(s\).
Write \(T_\delta^{(k)}\) for the spatial dilation obtained by replacing
\(|\xi|\) by \(2^{-k}|\xi|\) in the symbol of \(T_\delta\), with
\(s\in[1,2]\) unchanged. Choose smooth annular projections \(P_k\),
dilates of one fixed multiplier, which equal one on this spatial
frequency support. Lemma~\ref{lem:square-dilates} applied to \(P_kf\)
and the Littlewood--Paley inequality bound this shell, on all dyadic
radius intervals together, by
\[
 C\delta^{\alpha-1}
 \left(\delta^{1/2}+\delta+
        [1+\log(1/\delta)]^{1/2-1/p}\delta^{1/2-\sigma}\right)\|f\|_p .
\]
The sum over \(j\) converges because \(\alpha-1/2-\sigma>0\).

The multiplier left away from the boundary is smooth, supported in
\(\rho\le2\), and vanishes to order two at zero. Split it into annuli
\(\rho\asymp 2^\ell\), \(\ell\le2\). After extracting \(2^{2\ell}\)
and rescaling to the unit annulus, the multipliers have uniformly
bounded derivatives. Their kernels satisfy Lemma~\ref{lem:square-dilates}
at \(\delta=1\) with uniform constants (also for \(B_2,B_p\), by
the integrable kernel bound). For each \(\ell\) their dilates now have
frequency scale \(2^{k+\ell}\); the same Littlewood--Paley estimate
applies, and \(\sum_{\ell\le2}2^{2\ell}<\infty\).

These estimates apply first to finite sets of \(k\) and finite shell
sums. For Schwartz \(f\) the multipliers converge in the Fourier integral
for each fixed \(r\), by the local integrability of the boundary power.
Fatou's lemma and then monotone convergence give the full estimate.
The linear map \(f\mapsto (r\partial_r B_r^\alpha f)_{r>0}\) therefore
extends by density to \(L^p_x(L^2(\dd r/r))\), giving the asserted
extension of its pointwise norm.
\end{proof}
